\documentclass[aspectratio=169,10pt]{beamer}
\usepackage[T1]{fontenc}
\usepackage[brazil]{babel}
\usepackage{booktabs}
\usepackage{array}
\usepackage{tikz}
\usetikzlibrary{positioning,arrows.meta}
\usepackage{pgfplots}
\usepgfplotslibrary{groupplots}
\pgfplotsset{compat=1.18}
\usetheme{default}
\usecolortheme{beaver}
\setbeamertemplate{navigation symbols}{}
\setbeamertemplate{footline}[frame number]
\newcommand{\code}[1]{\texttt{#1}}
\newcommand{\NaR}{\mathrm{NaR}}
\IfFileExists{../results/spmv-macros.tex}{%
  % Generated by scripts/summarize_benchmark.py from spmv-errors.csv
\newcommand{\BenchMatrices}{358}
\newcommand{\BenchNnz}{3{,}054{,}096}
\newcommand{\BenchMaxNnz}{49{,}920}
\newcommand{\BenchMedianTEight}{-0.7}
\newcommand{\BenchConvMedianTEight}{-1.2}
\newcommand{\BenchBelowOneTEight}{87}
\newcommand{\BenchInfTEight}{0}
\newcommand{\BenchSatTEight}{0}
\newcommand{\BenchOverflowTEight}{0}
\newcommand{\BenchFlushTEight}{0}
\newcommand{\BenchAboveOneTEight}{13}
\newcommand{\BenchVecPerNnzTEight}{1.23}
\newcommand{\BenchMedianEFourMThree}{-1.2}
\newcommand{\BenchConvMedianEFourMThree}{-1.6}
\newcommand{\BenchBelowOneEFourMThree}{64}
\newcommand{\BenchInfEFourMThree}{34}
\newcommand{\BenchSatEFourMThree}{0}
\newcommand{\BenchOverflowEFourMThree}{120}
\newcommand{\BenchFlushEFourMThree}{173}
\newcommand{\BenchAboveOneEFourMThree}{36}
\newcommand{\BenchVecPerNnzEFourMThree}{1.80}
\newcommand{\BenchVecReductionEFourMThree}{32}
\newcommand{\BenchTotalReductionEFourMThree}{16}
\newcommand{\BenchMedianEFiveMTwo}{-1.0}
\newcommand{\BenchConvMedianEFiveMTwo}{-1.3}
\newcommand{\BenchBelowOneEFiveMTwo}{73}
\newcommand{\BenchInfEFiveMTwo}{25}
\newcommand{\BenchSatEFiveMTwo}{0}
\newcommand{\BenchOverflowEFiveMTwo}{86}
\newcommand{\BenchFlushEFiveMTwo}{138}
\newcommand{\BenchAboveOneEFiveMTwo}{27}
\newcommand{\BenchVecPerNnzEFiveMTwo}{1.80}
\newcommand{\BenchVecReductionEFiveMTwo}{32}
\newcommand{\BenchTotalReductionEFiveMTwo}{16}
\newcommand{\BenchMedianTSixteen}{-3.2}
\newcommand{\BenchConvMedianTSixteen}{-3.6}
\newcommand{\BenchBelowOneTSixteen}{97}
\newcommand{\BenchInfTSixteen}{0}
\newcommand{\BenchSatTSixteen}{0}
\newcommand{\BenchOverflowTSixteen}{0}
\newcommand{\BenchFlushTSixteen}{0}
\newcommand{\BenchAboveOneTSixteen}{3}
\newcommand{\BenchVecPerNnzTSixteen}{1.47}
\newcommand{\BenchMedianFPSixteen}{-3.5}
\newcommand{\BenchConvMedianFPSixteen}{-3.7}
\newcommand{\BenchBelowOneFPSixteen}{75}
\newcommand{\BenchInfFPSixteen}{25}
\newcommand{\BenchSatFPSixteen}{0}
\newcommand{\BenchOverflowFPSixteen}{86}
\newcommand{\BenchFlushFPSixteen}{102}
\newcommand{\BenchAboveOneFPSixteen}{25}
\newcommand{\BenchVecPerNnzFPSixteen}{1.84}
\newcommand{\BenchVecReductionFPSixteen}{20}
\newcommand{\BenchTotalReductionFPSixteen}{9}
\newcommand{\BenchMedianTThirtyTwo}{-8.2}
\newcommand{\BenchConvMedianTThirtyTwo}{-8.6}
\newcommand{\BenchBelowOneTThirtyTwo}{100}
\newcommand{\BenchInfTThirtyTwo}{0}
\newcommand{\BenchSatTThirtyTwo}{0}
\newcommand{\BenchOverflowTThirtyTwo}{0}
\newcommand{\BenchFlushTThirtyTwo}{0}
\newcommand{\BenchAboveOneTThirtyTwo}{0}
\newcommand{\BenchVecPerNnzTThirtyTwo}{2.14}
\newcommand{\BenchMedianFPThirtyTwo}{-7.6}
\newcommand{\BenchConvMedianFPThirtyTwo}{-7.8}
\newcommand{\BenchBelowOneFPThirtyTwo}{100}
\newcommand{\BenchInfFPThirtyTwo}{0}
\newcommand{\BenchSatFPThirtyTwo}{0}
\newcommand{\BenchOverflowFPThirtyTwo}{0}
\newcommand{\BenchFlushFPThirtyTwo}{2}
\newcommand{\BenchAboveOneFPThirtyTwo}{0}
\newcommand{\BenchVecPerNnzFPThirtyTwo}{2.76}
\newcommand{\BenchVecReductionFPThirtyTwo}{22}
\newcommand{\BenchTotalReductionFPThirtyTwo}{11}
\newcommand{\BenchMedianTEightStep}{-0.7}
\newcommand{\BenchConvMedianTEightStep}{-1.2}
\newcommand{\BenchBelowOneTEightStep}{87}
\newcommand{\BenchInfTEightStep}{0}
\newcommand{\BenchSatTEightStep}{0}
\newcommand{\BenchOverflowTEightStep}{0}
\newcommand{\BenchFlushTEightStep}{0}
\newcommand{\BenchAboveOneTEightStep}{13}
\newcommand{\BenchVecPerNnzTEightStep}{1.23}
\newcommand{\BenchMedianTSixteenStep}{-3.2}
\newcommand{\BenchConvMedianTSixteenStep}{-3.6}
\newcommand{\BenchBelowOneTSixteenStep}{97}
\newcommand{\BenchInfTSixteenStep}{0}
\newcommand{\BenchSatTSixteenStep}{0}
\newcommand{\BenchOverflowTSixteenStep}{0}
\newcommand{\BenchFlushTSixteenStep}{0}
\newcommand{\BenchAboveOneTSixteenStep}{3}
\newcommand{\BenchVecPerNnzTSixteenStep}{1.47}
\newcommand{\BenchMedianTThirtyTwoStep}{-8.2}
\newcommand{\BenchConvMedianTThirtyTwoStep}{-8.6}
\newcommand{\BenchBelowOneTThirtyTwoStep}{100}
\newcommand{\BenchInfTThirtyTwoStep}{0}
\newcommand{\BenchSatTThirtyTwoStep}{0}
\newcommand{\BenchOverflowTThirtyTwoStep}{0}
\newcommand{\BenchFlushTThirtyTwoStep}{0}
\newcommand{\BenchAboveOneTThirtyTwoStep}{0}
\newcommand{\BenchVecPerNnzTThirtyTwoStep}{2.14}
\newcommand{\BenchInRangeTEightEFourMThree}{235}
\newcommand{\BenchBetterTEightEFourMThree}{13}
\newcommand{\BenchWorseTEightEFourMThree}{181}
\newcommand{\BenchEqualTEightEFourMThree}{41}
\newcommand{\BenchInRangeTEightEFiveMTwo}{269}
\newcommand{\BenchBetterTEightEFiveMTwo}{45}
\newcommand{\BenchWorseTEightEFiveMTwo}{167}
\newcommand{\BenchEqualTEightEFiveMTwo}{57}
\newcommand{\BenchInRangeTSixteenFPSixteen}{270}
\newcommand{\BenchBetterTSixteenFPSixteen}{48}
\newcommand{\BenchWorseTSixteenFPSixteen}{110}
\newcommand{\BenchEqualTSixteenFPSixteen}{112}
\newcommand{\BenchInRangeTThirtyTwoFPThirtyTwo}{358}
\newcommand{\BenchBetterTThirtyTwoFPThirtyTwo}{199}
\newcommand{\BenchWorseTThirtyTwoFPThirtyTwo}{33}
\newcommand{\BenchEqualTThirtyTwoFPThirtyTwo}{126}
\newcommand{\BenchInRangeTEightTEightStep}{358}
\newcommand{\BenchBetterTEightTEightStep}{52}
\newcommand{\BenchWorseTEightTEightStep}{51}
\newcommand{\BenchEqualTEightTEightStep}{255}
\newcommand{\BenchInRangeTSixteenTSixteenStep}{358}
\newcommand{\BenchBetterTSixteenTSixteenStep}{10}
\newcommand{\BenchWorseTSixteenTSixteenStep}{13}
\newcommand{\BenchEqualTSixteenTSixteenStep}{335}
\newcommand{\BenchInRangeTThirtyTwoTThirtyTwoStep}{358}
\newcommand{\BenchBetterTThirtyTwoTThirtyTwoStep}{5}
\newcommand{\BenchWorseTThirtyTwoTThirtyTwoStep}{0}
\newcommand{\BenchEqualTThirtyTwoTThirtyTwoStep}{353}
\def\BenchCDF{../results/spmv-cdf.csv}\def\BenchDraft{}}{%
  % Generated by scripts/summarize_benchmark.py from synthetic-errors.csv
\newcommand{\BenchMatrices}{64}
\newcommand{\BenchNnz}{122{,}151}
\newcommand{\BenchMaxNnz}{6{,}712}
\newcommand{\BenchMedianTEight}{-0.6}
\newcommand{\BenchConvMedianTEight}{-0.8}
\newcommand{\BenchBelowOneTEight}{100}
\newcommand{\BenchInfTEight}{0}
\newcommand{\BenchSatTEight}{0}
\newcommand{\BenchOverflowTEight}{0}
\newcommand{\BenchFlushTEight}{0}
\newcommand{\BenchAboveOneTEight}{0}
\newcommand{\BenchVecPerNnzTEight}{0.98}
\newcommand{\BenchMedianEFourMThree}{\infty}
\newcommand{\BenchConvMedianEFourMThree}{-1.6}
\newcommand{\BenchBelowOneEFourMThree}{39}
\newcommand{\BenchInfEFourMThree}{61}
\newcommand{\BenchSatEFourMThree}{0}
\newcommand{\BenchOverflowEFourMThree}{30}
\newcommand{\BenchFlushEFourMThree}{57}
\newcommand{\BenchAboveOneEFourMThree}{61}
\newcommand{\BenchVecPerNnzEFourMThree}{1.42}
\newcommand{\BenchVecReductionEFourMThree}{32}
\newcommand{\BenchTotalReductionEFourMThree}{16}
\newcommand{\BenchMedianEFiveMTwo}{-1.1}
\newcommand{\BenchConvMedianEFiveMTwo}{-1.3}
\newcommand{\BenchBelowOneEFiveMTwo}{55}
\newcommand{\BenchInfEFiveMTwo}{45}
\newcommand{\BenchSatEFiveMTwo}{0}
\newcommand{\BenchOverflowEFiveMTwo}{28}
\newcommand{\BenchFlushEFiveMTwo}{35}
\newcommand{\BenchAboveOneEFiveMTwo}{45}
\newcommand{\BenchVecPerNnzEFiveMTwo}{1.42}
\newcommand{\BenchVecReductionEFiveMTwo}{32}
\newcommand{\BenchTotalReductionEFiveMTwo}{16}
\newcommand{\BenchMedianTSixteen}{-2.9}
\newcommand{\BenchConvMedianTSixteen}{-3.2}
\newcommand{\BenchBelowOneTSixteen}{100}
\newcommand{\BenchInfTSixteen}{0}
\newcommand{\BenchSatTSixteen}{0}
\newcommand{\BenchOverflowTSixteen}{0}
\newcommand{\BenchFlushTSixteen}{0}
\newcommand{\BenchAboveOneTSixteen}{0}
\newcommand{\BenchVecPerNnzTSixteen}{1.28}
\newcommand{\BenchMedianFPSixteen}{-3.5}
\newcommand{\BenchConvMedianFPSixteen}{-3.7}
\newcommand{\BenchBelowOneFPSixteen}{56}
\newcommand{\BenchInfFPSixteen}{44}
\newcommand{\BenchSatFPSixteen}{0}
\newcommand{\BenchOverflowFPSixteen}{19}
\newcommand{\BenchFlushFPSixteen}{21}
\newcommand{\BenchAboveOneFPSixteen}{44}
\newcommand{\BenchVecPerNnzFPSixteen}{1.62}
\newcommand{\BenchVecReductionFPSixteen}{21}
\newcommand{\BenchTotalReductionFPSixteen}{10}
\newcommand{\BenchMedianTThirtyTwo}{-7.8}
\newcommand{\BenchConvMedianTThirtyTwo}{-8.0}
\newcommand{\BenchBelowOneTThirtyTwo}{100}
\newcommand{\BenchInfTThirtyTwo}{0}
\newcommand{\BenchSatTThirtyTwo}{0}
\newcommand{\BenchOverflowTThirtyTwo}{0}
\newcommand{\BenchFlushTThirtyTwo}{0}
\newcommand{\BenchAboveOneTThirtyTwo}{0}
\newcommand{\BenchVecPerNnzTThirtyTwo}{1.92}
\newcommand{\BenchMedianFPThirtyTwo}{-7.5}
\newcommand{\BenchConvMedianFPThirtyTwo}{-7.6}
\newcommand{\BenchBelowOneFPThirtyTwo}{100}
\newcommand{\BenchInfFPThirtyTwo}{0}
\newcommand{\BenchSatFPThirtyTwo}{0}
\newcommand{\BenchOverflowFPThirtyTwo}{0}
\newcommand{\BenchFlushFPThirtyTwo}{0}
\newcommand{\BenchAboveOneFPThirtyTwo}{0}
\newcommand{\BenchVecPerNnzFPThirtyTwo}{2.49}
\newcommand{\BenchVecReductionFPThirtyTwo}{23}
\newcommand{\BenchTotalReductionFPThirtyTwo}{11}
\newcommand{\BenchMedianTEightStep}{-0.6}
\newcommand{\BenchConvMedianTEightStep}{-0.8}
\newcommand{\BenchBelowOneTEightStep}{100}
\newcommand{\BenchInfTEightStep}{0}
\newcommand{\BenchSatTEightStep}{0}
\newcommand{\BenchOverflowTEightStep}{0}
\newcommand{\BenchFlushTEightStep}{0}
\newcommand{\BenchAboveOneTEightStep}{0}
\newcommand{\BenchVecPerNnzTEightStep}{0.98}
\newcommand{\BenchMedianTSixteenStep}{-2.9}
\newcommand{\BenchConvMedianTSixteenStep}{-3.2}
\newcommand{\BenchBelowOneTSixteenStep}{100}
\newcommand{\BenchInfTSixteenStep}{0}
\newcommand{\BenchSatTSixteenStep}{0}
\newcommand{\BenchOverflowTSixteenStep}{0}
\newcommand{\BenchFlushTSixteenStep}{0}
\newcommand{\BenchAboveOneTSixteenStep}{0}
\newcommand{\BenchVecPerNnzTSixteenStep}{1.28}
\newcommand{\BenchMedianTThirtyTwoStep}{-7.8}
\newcommand{\BenchConvMedianTThirtyTwoStep}{-8.0}
\newcommand{\BenchBelowOneTThirtyTwoStep}{100}
\newcommand{\BenchInfTThirtyTwoStep}{0}
\newcommand{\BenchSatTThirtyTwoStep}{0}
\newcommand{\BenchOverflowTThirtyTwoStep}{0}
\newcommand{\BenchFlushTThirtyTwoStep}{0}
\newcommand{\BenchAboveOneTThirtyTwoStep}{0}
\newcommand{\BenchVecPerNnzTThirtyTwoStep}{1.92}
\newcommand{\BenchInRangeTEightEFourMThree}{25}
\newcommand{\BenchBetterTEightEFourMThree}{0}
\newcommand{\BenchWorseTEightEFourMThree}{25}
\newcommand{\BenchEqualTEightEFourMThree}{0}
\newcommand{\BenchInRangeTEightEFiveMTwo}{35}
\newcommand{\BenchBetterTEightEFiveMTwo}{9}
\newcommand{\BenchWorseTEightEFiveMTwo}{26}
\newcommand{\BenchEqualTEightEFiveMTwo}{0}
\newcommand{\BenchInRangeTSixteenFPSixteen}{36}
\newcommand{\BenchBetterTSixteenFPSixteen}{6}
\newcommand{\BenchWorseTSixteenFPSixteen}{30}
\newcommand{\BenchEqualTSixteenFPSixteen}{0}
\newcommand{\BenchInRangeTThirtyTwoFPThirtyTwo}{64}
\newcommand{\BenchBetterTThirtyTwoFPThirtyTwo}{45}
\newcommand{\BenchWorseTThirtyTwoFPThirtyTwo}{19}
\newcommand{\BenchEqualTThirtyTwoFPThirtyTwo}{0}
\newcommand{\BenchInRangeTEightTEightStep}{64}
\newcommand{\BenchBetterTEightTEightStep}{19}
\newcommand{\BenchWorseTEightTEightStep}{17}
\newcommand{\BenchEqualTEightTEightStep}{28}
\newcommand{\BenchInRangeTSixteenTSixteenStep}{64}
\newcommand{\BenchBetterTSixteenTSixteenStep}{0}
\newcommand{\BenchWorseTSixteenTSixteenStep}{0}
\newcommand{\BenchEqualTSixteenTSixteenStep}{64}
\newcommand{\BenchInRangeTThirtyTwoTThirtyTwoStep}{64}
\newcommand{\BenchBetterTThirtyTwoTThirtyTwoStep}{0}
\newcommand{\BenchWorseTThirtyTwoTThirtyTwoStep}{0}
\newcommand{\BenchEqualTThirtyTwoTThirtyTwoStep}{64}
\def\BenchCDF{../results/synthetic-cdf.csv}%
  \def\BenchDraft{\textcolor{red}{[rascunho: matrizes sintéticas]} }}

% Generated by scripts/run_mlp.py from mlp-errors.csv
\newcommand{\MlpTrials}{64}
\newcommand{\MlpRuns}{576}
\newcommand{\MlpMacs}{9{,}216}
\newcommand{\MlpMedianTEight}{-0.9}
\newcommand{\MlpVecPerMacTEight}{0.317}
\newcommand{\MlpUsableTEight}{7}
\newcommand{\MlpLostLowTEight}{19}
\newcommand{\MlpLostHighTEight}{47}
\newcommand{\MlpMedianEFourMThree}{-1.2}
\newcommand{\MlpVecPerMacEFourMThree}{0.567}
\newcommand{\MlpUsableEFourMThree}{4}
\newcommand{\MlpLostLowEFourMThree}{100}
\newcommand{\MlpLostHighEFourMThree}{100}
\newcommand{\MlpVecReductionEFourMThree}{44}
\newcommand{\MlpTotalReductionEFourMThree}{26}
\newcommand{\MlpMedianEFiveMTwo}{-0.9}
\newcommand{\MlpVecPerMacEFiveMTwo}{0.567}
\newcommand{\MlpUsableEFiveMTwo}{7}
\newcommand{\MlpLostLowEFiveMTwo}{2}
\newcommand{\MlpLostHighEFiveMTwo}{100}
\newcommand{\MlpVecReductionEFiveMTwo}{44}
\newcommand{\MlpTotalReductionEFiveMTwo}{26}
\newcommand{\MlpMedianTSixteen}{-3.3}
\newcommand{\MlpVecPerMacTSixteen}{0.572}
\newcommand{\MlpUsableTSixteen}{9}
\newcommand{\MlpLostLowTSixteen}{0}
\newcommand{\MlpLostHighTSixteen}{0}
\newcommand{\MlpMedianFPSixteen}{-3.2}
\newcommand{\MlpVecPerMacFPSixteen}{0.822}
\newcommand{\MlpUsableFPSixteen}{8}
\newcommand{\MlpLostLowFPSixteen}{0}
\newcommand{\MlpLostHighFPSixteen}{100}
\newcommand{\MlpVecReductionFPSixteen}{30}
\newcommand{\MlpTotalReductionFPSixteen}{16}
\newcommand{\MlpMedianTThirtyTwo}{-8.1}
\newcommand{\MlpVecPerMacTThirtyTwo}{1.082}
\newcommand{\MlpUsableTThirtyTwo}{9}
\newcommand{\MlpLostLowTThirtyTwo}{0}
\newcommand{\MlpLostHighTThirtyTwo}{0}
\newcommand{\MlpMedianFPThirtyTwo}{-7.2}
\newcommand{\MlpVecPerMacFPThirtyTwo}{1.582}
\newcommand{\MlpUsableFPThirtyTwo}{9}
\newcommand{\MlpLostLowFPThirtyTwo}{0}
\newcommand{\MlpLostHighFPThirtyTwo}{0}
\newcommand{\MlpVecReductionFPThirtyTwo}{32}
\newcommand{\MlpTotalReductionFPThirtyTwo}{17}
\newcommand{\MlpMedianTEightStep}{-0.9}
\newcommand{\MlpVecPerMacTEightStep}{0.317}
\newcommand{\MlpUsableTEightStep}{7}
\newcommand{\MlpLostLowTEightStep}{16}
\newcommand{\MlpLostHighTEightStep}{41}
\newcommand{\MlpMedianTSixteenStep}{-3.3}
\newcommand{\MlpVecPerMacTSixteenStep}{0.572}
\newcommand{\MlpUsableTSixteenStep}{9}
\newcommand{\MlpLostLowTSixteenStep}{0}
\newcommand{\MlpLostHighTSixteenStep}{0}
\newcommand{\MlpMedianTThirtyTwoStep}{-8.1}
\newcommand{\MlpVecPerMacTThirtyTwoStep}{1.082}
\newcommand{\MlpUsableTThirtyTwoStep}{9}
\newcommand{\MlpLostLowTThirtyTwoStep}{0}
\newcommand{\MlpLostHighTThirtyTwoStep}{0}
\newcommand{\MlpFusedBetterTEight}{80}
\newcommand{\MlpFusedWorseTEight}{83}
\newcommand{\MlpFusedRunsTEight}{576}
\newcommand{\MlpFusedBetterTSixteen}{0}
\newcommand{\MlpFusedWorseTSixteen}{0}
\newcommand{\MlpFusedRunsTSixteen}{576}
\newcommand{\MlpFusedBetterTThirtyTwo}{0}
\newcommand{\MlpFusedWorseTThirtyTwo}{0}
\newcommand{\MlpFusedRunsTThirtyTwo}{576}

% Generated by scripts/isa_analysis.py
\newcommand{\IsaZeroEFourMThree}{2}
\newcommand{\IsaSubEFourMThree}{14}
\newcommand{\IsaInfEFourMThree}{0}
\newcommand{\IsaNaNEFourMThree}{2}
\newcommand{\IsaSpecialEFourMThree}{7.0}
\newcommand{\IsaZeroEFiveMTwo}{2}
\newcommand{\IsaSubEFiveMTwo}{6}
\newcommand{\IsaInfEFiveMTwo}{2}
\newcommand{\IsaNaNEFiveMTwo}{6}
\newcommand{\IsaSpecialEFiveMTwo}{6.2}
\newcommand{\IsaZeroFPSixteen}{2}
\newcommand{\IsaSubFPSixteen}{2{,}046}
\newcommand{\IsaInfFPSixteen}{2}
\newcommand{\IsaNaNFPSixteen}{2{,}046}
\newcommand{\IsaSpecialFPSixteen}{6.2}
\newcommand{\IsaZeroBFSixteen}{2}
\newcommand{\IsaSubBFSixteen}{254}
\newcommand{\IsaInfBFSixteen}{2}
\newcommand{\IsaNaNBFSixteen}{254}
\newcommand{\IsaSpecialBFSixteen}{0.8}
\newcommand{\IsaZeroFPThirtyTwo}{2}
\newcommand{\IsaSubFPThirtyTwo}{16{,}777{,}214}
\newcommand{\IsaInfFPThirtyTwo}{2}
\newcommand{\IsaNaNFPThirtyTwo}{16{,}777{,}214}
\newcommand{\IsaSpecialFPThirtyTwo}{0.8}
\newcommand{\IsaZeroTEight}{1}
\newcommand{\IsaSubTEight}{0}
\newcommand{\IsaInfTEight}{0}
\newcommand{\IsaNaNTEight}{1}
\newcommand{\IsaSpecialTEight}{0.8}
\newcommand{\IsaZeroTSixteen}{1}
\newcommand{\IsaSubTSixteen}{0}
\newcommand{\IsaInfTSixteen}{0}
\newcommand{\IsaNaNTSixteen}{1}
\newcommand{\IsaSpecialTSixteen}{0.003}
\newcommand{\IsaZeroTThirtyTwo}{1}
\newcommand{\IsaSubTThirtyTwo}{0}
\newcommand{\IsaInfTThirtyTwo}{0}
\newcommand{\IsaNaNTThirtyTwo}{1}
\newcommand{\IsaSpecialTThirtyTwo}{0.000}


\title{Aritmética Takum na extensão vetorial RISC-V}
\subtitle{Instruções Zvtakum e avaliação em SpMV e camadas densas}
\author{Arthur dos Santos Fernandes \and Bruno Lima Soares \and Lucas Rodrigues Prosperi Ayala
  \and Pedro Henrique Campos Lima \and Pedro Lopes Ozelin}
\institute{Arquitetura e Organização de Computadores, DELT/UFMG, 2026/2}
\date{17/11/2026}

\begin{document}

\begin{frame}
\titlepage
\end{frame}

\begin{frame}{Motivação}
\begin{itemize}
\item Unidades vetoriais acumulam formatos de baixa precisão: binary16, bfloat16,
  OFP8 E4M3 e E5M2. Cada um traz subnormais, infinitos, NaNs, modos de arredondamento.
\item Mais formatos significam mais instruções de conversão e mais estado de controle.
\item \textbf{Takum} (Hunhold, 2024): precisão afunilada, faixa dinâmica quase
  independente do tamanho, um único decodificador para todas as larguras.
\item Hunhold (2025) mostrou ganhos no AVX10.2 e citou RISC-V como trabalho futuro.
\end{itemize}
\vspace{1em}
\textbf{Pergunta do projeto:} como fica o RVV com Takum como tipo de elemento,
em instruções, estado de controle e erro numérico?
\end{frame}

\begin{frame}{O formato Takum linear}
\centering
\begin{tikzpicture}[box/.style={draw, minimum height=8mm, align=center, font=\small}]
\node[box, minimum width=8mm, fill=red!15] (s) {S};
\node[box, minimum width=8mm, fill=orange!20, right=0mm of s] (d) {D};
\node[box, minimum width=20mm, fill=yellow!25, right=0mm of d] (r) {R (3 bits)};
\node[box, minimum width=28mm, fill=green!15, right=0mm of r] (c) {C ($r$ bits, $0\le r\le 7$)};
\node[box, minimum width=34mm, fill=blue!10, right=0mm of c] (f) {F ($p=n-5-r$ bits)};
\end{tikzpicture}
\vspace{1em}
\begin{itemize}
\item $r=R$ se $D=1$, $r=7-R$ se $D=0$; $c$ vem de $C$; $e=(1-2S)(c+S)$.
\item Valor: $x=[(1-3S)+f]\,2^e$. S, D, R e C cabem nos 12 bits mais significativos.
\item Especiais: só $0$ e $\NaR$. Negação por complemento de dois.
\item Fora da faixa: \textbf{saturação}; nunca vira $0$, $\pm\infty$ ou $\NaR$.
\end{itemize}
\end{frame}

\begin{frame}{Fase 1: modelo escalar exato}
\begin{columns}
\column{0.5\textwidth}
\begin{itemize}
\item Python, só biblioteca padrão, frações exatas.
\item Um arredondamento por operação (par mais próximo no fluxo de bits).
\item T8, T16, T32 e agora T64 (acumulador).
\end{itemize}
\column{0.5\textwidth}
\textbf{Validação}
\begin{itemize}
\item Oráculo independente: todos os $256^2$ pares T8.
\item Ida e volta exaustiva T8/T16.
\item 548.793 comparações com a libtakum: 0 divergências.
\end{itemize}
\end{columns}
\end{frame}

\begin{frame}{Decisões de projeto e alternativas}
\small
\begin{tabular}{>{\raggedright\arraybackslash}p{0.22\textwidth}>{\raggedright\arraybackslash}p{0.31\textwidth}>{\raggedright\arraybackslash}p{0.37\textwidth}}
\toprule
Decisão & Alternativas consideradas & Por quê\\
\midrule
Takum linear & logarítmico & Enunciado usa campo $F$; estudo do AVX10.2 usa linear\\
Simulador funcional em Python & RTL (VHDL/Verilog), C++ & Prazo; exatidão verificável; o enunciado aceita\\
Opcode custom-2 & funct6 livres do OP-V & Não colide com extensões padrão futuras\\
SEW seleciona T8/T16/T32 & bit novo em \code{vtype} & Nenhum estado novo\\
\code{vtdot}: soma exata, 1 arredondamento & arredondar a cada passo & Proposta de Hunhold; ablação mede o efeito\\
Acumulador $2\cdot$SEW em todos & larguras diferentes por formato & Comparação justa com o IEEE\\
\bottomrule
\end{tabular}
\end{frame}

\begin{frame}{Zvtakum: as cinco instruções}
\centering
\begin{tabular}{lll}
\toprule
Instrução & funct6 & Semântica\\
\midrule
\code{vtadd.vv} & \code{000000} & $v_d[i]=\mathrm{rd}(v_{s2}[i]+v_{s1}[i])$\\
\code{vtsub.vv} & \code{000010} & $v_d[i]=\mathrm{rd}(v_{s2}[i]-v_{s1}[i])$\\
\code{vtmul.vv} & \code{100100} & $v_d[i]=\mathrm{rd}(v_{s2}[i]\cdot v_{s1}[i])$\\
\code{vtdot.vv} & \code{111100} & $v_d[0]=\mathrm{rd}_{2\mathrm{SEW}}(v_d[0]+\sum v_{s2}[i]\,v_{s1}[i])$\\
\code{vncvt.t.t} & \code{010010} & $v_d[i]=\mathrm{rd}(v_{s2}[i])$, $2\mathrm{SEW}\to\mathrm{SEW}$\\
\bottomrule
\end{tabular}
\vspace{1em}
\begin{itemize}
\item Leiaute OP-V, opcode \textbf{custom-2}, funct3 OPFVV; funct6 igual ao do equivalente IEEE.
\item Máscara em \code{v0}; cauda e elementos inativos inalterados.
\item \code{vtdot}: T8$\to$T16, T16$\to$T32, T32$\to$T64.
\end{itemize}
\end{frame}

\begin{frame}{Nenhum CSR novo}
\centering
\begin{tabular}{ll}
\toprule
IEEE (\code{fcsr}) & Zvtakum\\
\midrule
\code{frm}: 5 modos de arredondamento & um único arredondamento\\
NV, DZ & $\NaR$ propagado no próprio dado\\
OF, UF & saturação: resultado finito e ordenado\\
NX & não sinalizado\\
subnormais, $\pm0$, $\pm\infty$, NaNs & não existem\\
E4M3 vs E5M2 no mesmo SEW & um formato por largura\\
\bottomrule
\end{tabular}
\end{frame}

\begin{frame}{Complexidade da ISA e do decodificador}
\begin{columns}[T]
\begin{column}{0.55\textwidth}
\footnotesize
\begin{tabular}{@{}>{\raggedright\arraybackslash}p{0.22\linewidth}>{\raggedright\arraybackslash}p{0.47\linewidth}>{\raggedright\arraybackslash}p{0.27\linewidth}@{}}
\toprule
Necessidade & RVV IEEE (+ BF16, OFP8) & Zvtakum\\
\midrule
Formato & SEW + \code{vtype.altfmt} & SEW\\
Ampliar & \code{vfwcvt}, \code{vfwcvtbf16} & \code{vwsll.vx}\\
Reduzir & \code{vfncvt[.rod]}, \code{vfncvtbf16[.sat]}, \code{vfncvt[.sat].f.f.q} & \code{vncvt.t.t}\\
Comparar, máx., ReLU & \code{vmf*}, \code{vfmin}, \code{vfmax} & \code{vms*}, \code{vmax} (inteiro)\\
Sinal, classe & \code{vfsgnj*}, \code{vfclass} & \code{vrsub}, \code{vmseq}\\
Estado & \code{frm}, \code{fflags} & nenhum\\
\bottomrule
\end{tabular}
\end{column}
\begin{column}{0.42\textwidth}
\footnotesize
\begin{tabular}{@{}lrrr@{}}
\toprule
Formato & Subnormais & NaN & Especiais\\
\midrule
E4M3 & \IsaSubEFourMThree & \IsaNaNEFourMThree & \IsaSpecialEFourMThree\%\\
E5M2 & \IsaSubEFiveMTwo & \IsaNaNEFiveMTwo & \IsaSpecialEFiveMTwo\%\\
binary16 & \IsaSubFPSixteen & \IsaNaNFPSixteen & \IsaSpecialFPSixteen\%\\
bfloat16 & \IsaSubBFSixteen & \IsaNaNBFSixteen & \IsaSpecialBFSixteen\%\\
\midrule
T$n$ & 0 & 1 (NaR) & $2/2^n$\\
\bottomrule
\end{tabular}
\vspace{0.8em}

Contagem exaustiva de todos os padrões (\code{make isa}). Takum: uma regra
para toda largura, expoente nos 12 primeiros bits.
\end{column}
\end{columns}
\vspace{0.8em}
\small Takum é ordenado como inteiro em complemento de dois: comparação, ReLU,
negação e ampliação usam instruções de inteiro que o RVV já tem.
\end{frame}

\begin{frame}[fragile]{Kernel SpMV CSR (laço interno, por faixa)}
\footnotesize
\begin{columns}[T]
\column{0.32\textwidth}
\textbf{Takum: 5}
\begin{verbatim}
vsetvli  e(SEW)
vle      valores
vle32    colunas
vluxei32 x[col]
vtdot.vv
\end{verbatim}
\column{0.32\textwidth}
\textbf{binary16/32: 7}
\begin{verbatim}
vsetvli  e(SEW)
vle / vle32 / vluxei32
vfwmul.vv
vsetvli  e(2SEW)
vfredosum.vs
\end{verbatim}
\column{0.32\textwidth}
\textbf{FP8: 9}
\begin{verbatim}
vsetvli  e8
vle / vle32 / vluxei32
vfwcvt x2  (FP8->FP16)
vsetvli  e16
vfmul.vv
vfredosum.vs
\end{verbatim}
\end{columns}
\vspace{0.5em}
\normalsize
Mesmo controle escalar em todos (6 por faixa, 7 por linha). O RVV ratificado não
tem FP8: as conversões FP8$\leftrightarrow$FP16 são modeladas.
\end{frame}

\begin{frame}{Metodologia da avaliação}
\begin{itemize}
\item \textbf{Matrizes:} \BenchMatrices{} matrizes reais do NIST Matrix Market com até
  50 mil não nulos (a SuiteSparse recusou os downloads); valores decimais lidos exatamente.
\item \textbf{Formatos:} T8, E4M3, E5M2 · T16, binary16 · T32, binary32.
\item \textbf{Referência:} binary128 (entradas, produtos e somas); $x$ = uns.
\item \textbf{Erro:} $\|y_f-y_{\mathrm{ref}}\|_2/\|y_{\mathrm{ref}}\|_2$, calculado
  exatamente; NaN ou $\infty$ conta como erro infinito.
\item \textbf{Instruções:} VLEN = 128, LMUL = 1.
\item \textbf{Confiabilidade:} o caminho rápido do benchmark é idêntico, bit a bit, ao
  simulador instrução por instrução (testado), e a fórmula de contagem é igual à do simulador.
\end{itemize}
\end{frame}

\begin{frame}{Resultado: distribuição do erro \BenchDraft}
\centering
\begin{tikzpicture}
\begin{groupplot}[group style={group size=3 by 1, horizontal sep=12mm},
  width=0.36\textwidth, height=5cm, xmin=0, xmax=100, grid=major, grid style={gray!25},
  tick label style={font=\scriptsize}, label style={font=\scriptsize},
  title style={font=\small}, legend style={font=\scriptsize, at={(0.02,0.96)}, anchor=north west},
  unbounded coords=jump, no markers, every axis plot/.append style={thick},
  xlabel={matrizes (\%)}]
\nextgroupplot[title={8 bits}, ymin=-4, ymax=2, y filter/.expression={y<-4 ? -4 : y}, ylabel={$\log_{10}$ erro relativo}]
\addplot table[x=share, y=T8, col sep=comma] {\BenchCDF}; \addlegendentry{T8}
\addplot table[x=share, y=E4M3, col sep=comma] {\BenchCDF}; \addlegendentry{E4M3}
\addplot table[x=share, y=E5M2, col sep=comma] {\BenchCDF}; \addlegendentry{E5M2}
\nextgroupplot[title={16 bits}, ymin=-8, ymax=2, y filter/.expression={y<-8 ? -8 : y}]
\addplot table[x=share, y=T16, col sep=comma] {\BenchCDF}; \addlegendentry{T16}
\addplot table[x=share, y=FP16, col sep=comma] {\BenchCDF}; \addlegendentry{binary16}
\nextgroupplot[title={32 bits}, ymin=-14, ymax=2, y filter/.expression={y<-14 ? -14 : y}]
\addplot table[x=share, y=T32, col sep=comma] {\BenchCDF}; \addlegendentry{T32}
\addplot table[x=share, y=FP32, col sep=comma] {\BenchCDF}; \addlegendentry{binary32}
\end{groupplot}
\end{tikzpicture}

\small A curva termina onde as matrizes restantes deram NaN ou $\infty$.
\end{frame}

\begin{frame}{Resultado: resumo \BenchDraft}
\centering
\begin{tabular}{lrrrr}
\toprule
Formato & Mediana $\log_{10}$ & erro $<100\%$ & NaN/$\infty$ & instr. vetoriais/nnz\\
\midrule
T8 & $\BenchMedianTEight$ & \BenchBelowOneTEight\% & \BenchInfTEight\% & \BenchVecPerNnzTEight\\
E4M3 & $\BenchMedianEFourMThree$ & \BenchBelowOneEFourMThree\% & \BenchInfEFourMThree\% & \BenchVecPerNnzEFourMThree\\
E5M2 & $\BenchMedianEFiveMTwo$ & \BenchBelowOneEFiveMTwo\% & \BenchInfEFiveMTwo\% & \BenchVecPerNnzEFiveMTwo\\
\midrule
T16 & $\BenchMedianTSixteen$ & \BenchBelowOneTSixteen\% & \BenchInfTSixteen\% & \BenchVecPerNnzTSixteen\\
binary16 & $\BenchMedianFPSixteen$ & \BenchBelowOneFPSixteen\% & \BenchInfFPSixteen\% & \BenchVecPerNnzFPSixteen\\
\midrule
T32 & $\BenchMedianTThirtyTwo$ & \BenchBelowOneTThirtyTwo\% & \BenchInfTThirtyTwo\% & \BenchVecPerNnzTThirtyTwo\\
binary32 & $\BenchMedianFPThirtyTwo$ & \BenchBelowOneFPThirtyTwo\% & \BenchInfFPThirtyTwo\% & \BenchVecPerNnzFPThirtyTwo\\
\bottomrule
\end{tabular}
\vspace{1em}

Instruções vetoriais a menos com Takum: \BenchVecReductionEFourMThree\% contra FP8,
\BenchVecReductionFPSixteen\% contra binary16, \BenchVecReductionFPThirtyTwo\% contra binary32.
\end{frame}

\begin{frame}{Aprendizado de máquina: camada densa $y=\mathrm{ReLU}(Wx)$}
\begin{columns}[T]
\begin{column}{0.58\textwidth}
\begin{tikzpicture}
\begin{groupplot}[group style={group size=1 by 3, vertical sep=2mm, xticklabels at=edge bottom},
  width=0.8\linewidth, height=2.6cm, xmin=-16, xmax=16, xtick={-16,-8,...,16},
  grid=major, grid style={gray!25}, tick label style={font=\tiny}, label style={font=\scriptsize},
  legend style={font=\tiny, at={(1.02,0.5)}, anchor=west, cells={anchor=west}},
  every axis plot/.append style={thick, mark size=1pt}]
\nextgroupplot[ymin=-1.5, ymax=0.3]
\addplot table[x=scale, y=T8, col sep=comma] {../results/mlp-curve.csv}; \addlegendentry{T8}
\addplot table[x=scale, y=E4M3, col sep=comma] {../results/mlp-curve.csv}; \addlegendentry{E4M3}
\addplot table[x=scale, y=E5M2, col sep=comma] {../results/mlp-curve.csv}; \addlegendentry{E5M2}
\nextgroupplot[ymin=-3.6, ymax=0.3, ylabel={mediana $\log_{10}$ erro}]
\addplot table[x=scale, y=T16, col sep=comma] {../results/mlp-curve.csv}; \addlegendentry{T16}
\addplot table[x=scale, y=FP16, col sep=comma] {../results/mlp-curve.csv}; \addlegendentry{binary16}
\nextgroupplot[ymin=-8.5, ymax=-6.8, xlabel={escala da entrada $2^k$}]
\addplot table[x=scale, y=T32, col sep=comma] {../results/mlp-curve.csv}; \addlegendentry{T32}
\addplot table[x=scale, y=FP32, col sep=comma] {../results/mlp-curve.csv}; \addlegendentry{binary32}
\end{groupplot}
\end{tikzpicture}
\end{column}
\begin{column}{0.4\textwidth}
\small
\begin{itemize}
\item \MlpTrials{} redes aleatórias $64\to64\to64\to16$, sem bias; entrada $\times 2^k$.
\item ReLU Takum: \code{vmax.vx} de inteiro.
\item IEEE: erro constante na faixa, perde a saída fora dela (E4M3 em
  \the\numexpr 9-\MlpUsableEFourMThree\relax{} de 9 escalas).
\item T16 e T32 nunca perdem; T8 degrada aos poucos.
\item Em $k=0$: E4M3 $\MlpMedianEFourMThree$ contra T8 $\MlpMedianTEight$.
\item Instruções vetoriais: $-\MlpVecReductionEFourMThree\%$ (FP8),
  $-\MlpVecReductionFPSixteen\%$ (binary16).
\end{itemize}
\end{column}
\end{columns}
\end{frame}

\begin{frame}{Discussão e limitações}
\textbf{Precisão dentro da faixa:} onde o IEEE não estoura, ele costuma ganhar em
8 e 16 bits (T8 perde para E4M3 em \BenchWorseTEightEFourMThree{} de
\BenchInRangeTEightEFourMThree{} matrizes; T16 perde para binary16 em
\BenchWorseTSixteenFPSixteen{} de \BenchInRangeTSixteenFPSixteen). É o preço da
precisão afunilada: bits gastos em faixa dinâmica.
\vspace{0.6em}

\textbf{Soma exata no \code{vtdot}:} quase neutra para o erro com $x$ = uns. Mudou o
resultado em \the\numexpr\BenchBetterTEightTEightStep+\BenchWorseTEightTEightStep\relax{}
matrizes T8 (melhor em \BenchBetterTEightTEightStep, pior em \BenchWorseTEightTEightStep).
O ganho do \code{vtdot} aqui é em instruções.
\vspace{1em}

\textbf{Limitações}
\begin{itemize}
\item Modelo funcional: conta instruções, não ciclos, área nem energia.
\item Menos instruções $\neq$ mais rápido: \code{vtdot} é uma operação longa em hardware.
\item Linha de base FP8 usa conversões que o RVV ratificado não tem.
\item SpMV só com $x$ = uns; redes aleatórias, não treinadas.
\end{itemize}
\end{frame}

\begin{frame}[fragile]{Demonstração}
\begin{verbatim}
python3 -m takum vector-demo   # 5 instruções, rastro, codificação
python3 -m takum add --width 8 1.5 2.25
make test                      # suíte completa de testes
make bench                     # benchmark, retomável
make mlp                       # camadas densas, varredura de escala
make isa                       # contagem de padrões especiais
make paper slides
\end{verbatim}
\end{frame}

\begin{frame}{Conclusão}
\begin{itemize}
\item Takum entra no RVV como formato selecionado por SEW: cinco instruções,
  opcode custom, \textbf{nenhum estado de controle novo}.
\item Em 8 e 16 bits, Takum sempre deu resultado finito; OFP8 e binary16 falharam em
  de \BenchInfFPSixteen\% a \BenchInfEFourMThree\% das matrizes, mas são mais precisos onde não estouram.
\item Em 32 bits, T32 foi mais preciso que binary32 (\BenchBetterTThirtyTwoFPThirtyTwo{} contra \BenchWorseTThirtyTwoFPThirtyTwo{} matrizes).
\item Camadas densas: T16 e T32 mantêm a saída em toda a faixa $2^{\pm16}$; E4M3 precisa de fatores de escala.
\item O \code{vtdot} encurta os laços da SpMV e da camada densa; comparação, ReLU e ampliação usam instruções de inteiro.
\item Próximos passos: modelo de ciclos ou RTL do \code{vtdot}, Takum logarítmico,
  redes treinadas.
\end{itemize}
\vspace{1em}
\centering \small \url{https://github.com/pedrohcali/aoc-takum-rvv}
\end{frame}

\end{document}
